\begin{scanpage}{0201}{183}
\sectiontitle{29}{Radiation of a Transmitting Vibrator}
The radiation characteristic of a transmitting vibrator may be calculated from a known equation by integrating the currents in it. However, such an approach is not advisable because, as was indicated above, the precision of Equations (28.03) is different in different parts of the conductor, and is low near its ends ($z=z_1$ and $z=z_2$) and the point $z=0$. The principle of duality gives more precise results. This principle leads to the following expression for the radiation field in the far zone [82]:
\[
\left.\begin{aligned}
E_\vartheta=H_\varphi&=\frac{\mathcal E}{2\sin\vartheta\ln\dfrac{2i}{\gamma ka\sin\vartheta}}\frac{e^{ikR}}R\cdot f(\vartheta),\\
E_\varphi=H_\vartheta&=0.
\end{aligned}\right\}\tag{29.01}
\]
The function
\[
\begin{aligned}
f(\vartheta)={}&1-\psi_+(z_2)e^{ikz_2(1-\cos\vartheta)}-\psi_-(-z_1)e^{-ikz_1(1+\cos\vartheta)}\\
&+B_1\psi_+(L)e^{ikz_2(1-\cos\vartheta)}+B_2\psi_-(L)e^{-ikz_1(1+\cos\vartheta)}
\end{aligned}\tag{29.02}
\]
is connected with the current (28.14) excited in a vibrator by plane wave (28.13) by the relationship
\[J(0)=Sf(\vartheta).\tag{29.03}\]
The coefficients $B_1$ and $B_2$ do not depend on the angle $\vartheta$.

Expression (29.02) enables one to trace the formation of the radiation. The first term (one) is the radiation field of an infinitely long conductor excited by a concentrated emf. Propagating in the direction $\vartheta=0$, this field reaches the conductor's end $z=z_2$ and---being diffracted by it---generates a primary edge wave (the second term). In a similar way, the primary edge wave diverging from the conductor's end $z=z_1$ is excited (the third term). The last two terms in Equation (29.02) determine the waves arising as a result of subsequent diffraction (secondary, tertiary, etc.). The amplitudes $B_1$ and $B_2$ of these waves may be found from the conditions
\end{scanpage}
